% Gesamttest «Planimetrie» — Quelle des PDFs (python3 scripts/build-lp-pdf.py planimetrie).
% Musterlösung und Raster: bewertungspaket.tex. Alle Werte mit python3 nachgerechnet (06.10.2026).
% Kein \lt/\gt — pdfLaTeX kennt sie nicht, < und > tun es.
\documentclass[11pt]{article}
\usepackage{../lp-druck}
\renewcommand{\lpname}{Planimetrie}
\renewcommand{\lpteilgebiet}{5.2}
\renewcommand{\lpbereich}{Grundlagenfach}
\renewcommand{\lpfuss}{Gesamttest · 25 Punkte}
\begin{document}
\lpkopf{Gesamttest}{%
  \vspace{4pt}\noindent\begin{tabularx}{\linewidth}{@{}X X@{}}
  Code (kein Name): \hrulefill & Punkte: \hrulefill\ / 25
  \end{tabularx}}

\begin{lpinfo}{Anleitung}
\textbf{Zeit:} rund 30 Minuten. \textbf{Hilfsmittel:} Taschenrechner erlaubt (RLP 5.2 trägt keinen Vermerk «auch ohne Hilfsmittel»).
Mach zu jeder Aufgabe eine \textbf{Skizze} und schreib den \textbf{Rechenweg} auf — bewertet wird der Weg. Ergebnisse mit
\textbf{Einheit}, gerundet auf zwei Dezimalen. Die Figuren sind nicht massstäblich.
Danach: Lösung scannen oder fotografieren (ohne Namen und Standort) und mit dem \emph{Bewertungspaket} bewerten lassen.
\end{lpinfo}

\lpteil{Teil A · Dreiecke}{Kapitel 1 und 2 · 7 Punkte}

\begin{lpaufgabe}{G1}{Dreiecke beschreiben}{3 P}
Ein gleichschenkliges Dreieck hat die Basiswinkel $64^\circ$.
\begin{enumerate}[label=(\alph*),leftmargin=*,itemsep=1pt]
\item Wie gross ist der Spitzenwinkel? Begründe.
\item Was ist die Seitenhalbierende $s_a$? Beschreib sie in einem Satz.
\end{enumerate}
\lpplatz{4}
\end{lpaufgabe}

\begin{lpaufgabe}{G2}{Dreiecksfläche und Höhe}{4 P}
\begin{minipage}[t]{0.52\linewidth}\vspace{0pt}
Im Dreieck $ABC$ ist die Grundseite $c = \overline{AB} = 9\,\text{cm}$ und die zugehörige Höhe $h_c = 4.2\,\text{cm}$.
Die Seite $b = \overline{AC}$ ist $7\,\text{cm}$ lang.
\begin{enumerate}[label=(\alph*),leftmargin=*,itemsep=1pt]
\item Berechne den Flächeninhalt $A$.
\item Berechne die Höhe $h_b$ zur Seite $b$.
\end{enumerate}
\end{minipage}\hfill
\begin{minipage}[t]{0.44\linewidth}\vspace{0pt}\centering
\begin{tikzpicture}[scale=0.42]
  \coordinate (A) at (0,0); \coordinate (B) at (9,0); \coordinate (C) at (-2.2,4.2);
  \draw[lpblau,thick,fill=lpblau!8] (A)--(B)--(C)--cycle;
  \draw[lpgrau,dashed] (A)--(-2.2,0);
  \draw[lporange,thick,dashed] (C)--(-2.2,0) node[midway,left] {\small $h_c$};
  \draw[lporange] (-2.2,0) rectangle (-1.8,0.4);
  \node[below] at (A) {$A$}; \node[below] at (B) {$B$}; \node[above] at (C) {$C$};
  \node[below] at (4.5,0) {\small $c = 9$};
  \node[right] at (-0.9,2.2) {\small $b$};
\end{tikzpicture}\\[-2pt]{\footnotesize nicht massstäblich}
\end{minipage}
\lpplatz{4}
\end{lpaufgabe}

\lpteil{Teil B · Vierecke und Kreis}{Kapitel 3 und 4 · 12 Punkte}

\begin{lpaufgabe}{G3}{Gleichschenkliges Trapez}{5 P}
\begin{minipage}[t]{0.52\linewidth}\vspace{0pt}
Ein gleichschenkliges Trapez hat die Parallelseiten $a = 11\,\text{cm}$ und $c = 5\,\text{cm}$ und die Höhe $h = 6\,\text{cm}$.
Berechne den Flächeninhalt, die Länge der Mittellinie $m$, die Länge eines Schenkels und den Umfang.
\end{minipage}\hfill
\begin{minipage}[t]{0.44\linewidth}\vspace{0pt}\centering
\begin{tikzpicture}[scale=0.4]
  \coordinate (A) at (0,0); \coordinate (B) at (11,0); \coordinate (C) at (8,6); \coordinate (D) at (3,6);
  \draw[lpblau,thick,fill=lpblau!8] (A)--(B)--(C)--(D)--cycle;
  \draw[lporange,dashed] (D)--(3,0) node[midway,right] {\small $h$};
  \node[below] at (5.5,0) {\small $a = 11$}; \node[above] at (5.5,6) {\small $c = 5$};
  \node[below left] at (A) {$A$}; \node[below right] at (B) {$B$}; \node[above right] at (C) {$C$}; \node[above left] at (D) {$D$};
\end{tikzpicture}
\end{minipage}
\lpplatz{6}
\end{lpaufgabe}

\begin{lpaufgabe}{G4}{Kreissektor}{4 P}
\begin{minipage}[t]{0.6\linewidth}\vspace{0pt}
Ein Kreissektor hat den Radius $r = 8\,\text{cm}$ und den Mittelpunktswinkel $\varphi = 135^\circ$.
Berechne die Bogenlänge $b$ und den Flächeninhalt $A_S$ des Sektors.
\end{minipage}\hfill
\begin{minipage}[t]{0.36\linewidth}\vspace{0pt}\centering
\begin{tikzpicture}[scale=0.28]
  \draw[lpgrau] (0,0) circle (8);
  \fill[lpgruen!20] (0,0) -- (8,0) arc (0:135:8) -- cycle;
  \draw[lpgruen,thick] (0,0) -- (8,0) arc (0:135:8) -- cycle;
  \draw[lporange] (2,0) arc (0:135:2); \node at (1.2,2.4) {\small $\varphi$};
  \node[below] at (4,0) {\small $r$};
\end{tikzpicture}
\end{minipage}
\lpplatz{5}
\end{lpaufgabe}

\begin{lpaufgabe}{G5}{Kreissegment}{3 P}
\begin{minipage}[t]{0.6\linewidth}\vspace{0pt}
Eine Sehne schneidet von einem Kreis mit $r = 10\,\text{cm}$ ein Segment ab; der zugehörige Mittelpunktswinkel ist $\varphi = 60^\circ$.
Berechne den Flächeninhalt des Segments (grün).
\end{minipage}\hfill
\begin{minipage}[t]{0.36\linewidth}\vspace{0pt}\centering
\begin{tikzpicture}[scale=0.24]
  \draw[lpgrau] (0,0) circle (10);
  \fill[lpgruen!25] (10,0) arc (0:60:10) -- cycle;
  \draw[lpblau,thick] (0,0) -- (10,0) -- (5,8.66) -- cycle;
  \draw[lpgruen,thick] (10,0) arc (0:60:10);
  \draw[lporange] (2.5,0) arc (0:60:2.5); \node at (3.2,1.6) {\small $\varphi$};
\end{tikzpicture}
\end{minipage}
\lpplatz{5}
\end{lpaufgabe}

\lpteil{Teil C · Ähnlichkeit}{Kapitel 5 · 6 Punkte}

\begin{lpaufgabe}{G6}{Schatten}{3 P}
Eine 1.8\,m grosse Person wirft einen 2.4\,m langen Schatten. Zur gleichen Zeit wirft ein Fahnenmast einen 14\,m langen Schatten.
Wie hoch ist der Mast? Mach eine Skizze mit den beiden ähnlichen Dreiecken.
\lpplatz{6}
\end{lpaufgabe}

\begin{lpaufgabe}{G7}{Massstab und Fläche}{3 P}
Ein Architekturmodell ist im Massstab $1 : 50$ gebaut. Die Grundfläche eines Raums misst im Modell $120\,\text{cm}^2$.
Wie gross ist sie in Wirklichkeit? Gib das Ergebnis in $\text{m}^2$ an.
\lpplatz{5}
\end{lpaufgabe}
\end{document}
